\chapter*{Anexo F: Repositorio del proyecto y guía de uso}
\addcontentsline{toc}{chapter}{Anexo F: Repositorio del proyecto y guía de uso}
\label{cap:anexo_repositorio}

El código fuente completo de la plataforma, junto con este documento, está
disponible públicamente en el siguiente repositorio, que es el mismo que
contiene el Trabajo Fin de Grado sobre el que se construye este trabajo:

\begin{center}
\url{https://github.com/holt00/TFG-open-cvn-schema}
\end{center}

Este anexo describe la organización del repositorio y orienta a quien quiera
consultarlo o ejecutarlo, sin repetir los comandos del
Anexo~\ref{cap:anexo_reproduccion}.

\section*{Estructura del repositorio}

\begin{itemize}
  \sloppy
  \item \path{src/tfm_lakehouse/}: código propio de la plataforma, con un
  subpaquete por etapa. \path{orcid_client} y \path{orcid_bulk} acceden a las
  dos fuentes de ORCID, \path{synthetic_cvn} genera los currículos
  sintéticos, \path{bronze} aterriza los registros, \path{silver} valida,
  normaliza y resuelve entidades, \path{gold} calcula los indicadores y
  genera el SQL de publicación, \path{spark_jobs} contiene los puntos de
  entrada de los trabajos de Spark y \path{benchmark} contiene el banco de
  pruebas.
  \item \path{dags/}: definición de los flujos \texttt{ingest\_validate} y
  \texttt{transform\_publish} de Airflow.
  \item \path{infra/}: todo lo necesario para desplegar la plataforma. Los
  valores de Helm de los servicios están en \path{helm-values/}, las
  imágenes y la configuración de Iceberg y de Spark en \path{spark-conf/} e
  \path{ingest/}, y el panel de Superset, su rol de solo lectura y su
  importación en \path{superset/}. Cada subdirectorio tiene su propio
  \path{README.md}.
  \item \path{tests/}: batería automatizada, que incluye las pruebas de los
  trabajos de Spark ejecutadas dentro de la imagen del clúster.
  \item \path{docs/benchmark/}: resultados completos del banco de pruebas, con
  los registros de cada ejecución en formato CSV y JSON, las gráficas y el
  procedimiento para repetir la campaña.
  \item \path{docs/roadmap/tfm/}: un documento por tarea, con su plan, las
  decisiones tomadas y las desviaciones respecto a lo planeado.
  \item \path{docs/pipeline/known_limitations.md}: registro de limitaciones
  resumido en el Anexo~\ref{cap:anexo_limitaciones}.
  \item \path{docs/memoria/TFM/}: código fuente de este documento.
  \item \path{src/open_cvn/}, \path{src/open_cvn_app/} y \path{schemas/}:
  formato Open CVN, herramienta local y esquema JSON del trabajo anterior, de
  los que la plataforma toma el contrato de validación del currículo.
\end{itemize}

\section*{Cómo orientarse en el repositorio}

El punto de entrada es el fichero \path{README.md} de la raíz, que resume el
objetivo y el alcance. El fichero \path{PROJECT_GUIDE.md} añade un mapa de la
documentación bajo \path{docs/}. Para estudiar el desarrollo en orden
cronológico sirve el registro de estado en
\path{docs/context/tfm/current_status.md}, que recoge cada tarea con sus
cifras medidas, y para entender una decisión concreta conviene abrir el
documento de la tarea correspondiente, cuyas tablas de decisiones enumeran
cada alternativa descartada con su motivo.

Para reproducir la plataforma, la guía mantenida es
\path{docs/development/tfm_lakehouse_workflow.md}, que el
Anexo~\ref{cap:anexo_reproduccion} resume. Para repetir solo la campaña de
rendimiento, el procedimiento está en \path{docs/benchmark/README.md}, y
trabaja sobre espacios de nombres aislados que no alteran las tablas de
producción.
